\documentclass[11pt,a4paper]{article}
\usepackage[margin=1in]{geometry}
\usepackage{amsmath,amssymb,amsthm}
\usepackage{algorithm}
\usepackage{algpseudocode}
\usepackage{booktabs}
\usepackage{hyperref}

\theoremstyle{definition}
\newtheorem{definition}{\textbf{Definition}}[section]
\newtheorem{theorem}{\textbf{Theorem}}[section]
\newtheorem{lemma}[theorem]{\textbf{Lemma}}
\newtheorem{remark}{\textbf{Remark}}[section]

\title{\textbf{Flow Matching and Rectified Flow}}
\author{\textbf{Team Scaibu} \\ \small Email: \texttt{jaydeep.9664920749@gmail.com} \\ \small \textbf{Architecture and Core Engine Team}}
\date{\textbf{October 2026}}

\begin{document}
\maketitle

\section{Definitions and Cognitive Mental Models}

\subsection{Formal Mathematical Definition}
\textbf{Definition (Conditional Flow Matching and Straight-Line Transport):}
Let $p_0$ denote the base Gaussian noise prior $\mathcal{N}(\mathbf{0}, \mathbf{I})$ and $p_1$ denote the empirical data distribution on $\mathbb{R}^D$. A Continuous Normalizing Flow defines a time-dependent probability path $p_t$ for $t \in [0, 1]$ generated by an Ordinary Differential Equation (ODE):
{\boldmath
\begin{equation}
\frac{\mathrm{d}\mathbf{x}}{\mathrm{d}t} = \mathbf{v}_t(\mathbf{x}), \quad \mathbf{x}(0) \sim p_0, \quad \mathbf{x}(1) \sim p_1
\end{equation}
}
Under the canonical \textbf{Rectified Flow / Optimal Transport Flow Matching} probability path, intermediate samples are defined by straight-line affine interpolation between noise $\mathbf{x}_0 \sim p_0$ and data $\mathbf{x}_1 \sim p_1$:
{\boldmath
\begin{equation}
\mathbf{x}_t = (1 - t)\mathbf{x}_0 + t \mathbf{x}_1, \quad t \in [0, 1]
\end{equation}
}
The conditional velocity vector field governing this path is constant across time:
{\boldmath
\begin{equation}
\mathbf{u}_t(\mathbf{x}_t \mid \mathbf{x}_0, \mathbf{x}_1) = \frac{\mathrm{d}\mathbf{x}_t}{\mathrm{d}t} = \mathbf{x}_1 - \mathbf{x}_0
\end{equation}
}
The parameterized neural vector field $\mathbf{v}_\theta(\mathbf{x}, t)$ is trained via the simulation-free \textbf{Conditional Flow Matching (CFM)} objective:
{\boldmath
\begin{equation}
\mathcal{L}_{\text{CFM}}(\theta) = \mathbb{E}_{t \sim \mathcal{U}(0, 1), \mathbf{x}_0 \sim p_0, \mathbf{x}_1 \sim p_1}\left[ \left\| \mathbf{v}_\theta(\mathbf{x}_t, t) - (\mathbf{x}_1 - \mathbf{x}_0) \right\|_2^2 \right]
\end{equation}
}
Because straight lines minimize geodesic transport cost, the resulting vector field has minimal curvature, enabling fast, high-quality generation in 1 to 4 integration steps.

\subsection{Expanded Non-Technical Definition (Intuition for Anyone)}
\textbf{Intuitive Conceptual Definition:}
\begin{enumerate}
    \item \textbf{Curved Highways vs Straight Flight Paths}: Traditional diffusion models move along curved, winding spirals through noise space. Because the path is curved, taking large jumps causes errors, forcing you to take dozens of tiny steps.
    \item \textbf{The Straightest Possible Line}: Flow Matching asks a simple question: ``Why take a winding path? Why not draw a perfectly straight line from random noise $\mathbf{x}_0$ directly to the target image $\mathbf{x}_1$?''
    \item \textbf{The Constant Speed Vector}: Along a straight line, your velocity is just $\mathbf{x}_1 - \mathbf{x}_0$. If you are at point $\mathbf{x}_t$ at time $t$, you just keep driving in that exact direction.
    \item \textbf{One-Step Teleportation}: Because the path is nearly straight, you don't need a complex multi-step ODE solver. You can compute the velocity vector once and take a single giant stride from $t = 0$ all the way to $t = 1$, generating an image instantly!
\end{enumerate}

\subsection{Child-Friendly Mental Model (Physical Metaphor)}
Imagine walking from your bedroom to the kitchen:
\begin{itemize}
    \item \textbf{Diffusion Process}: You wander through every room in the house, spinning in circles, bumping into walls, and slowly spiraling into the kitchen over 50 steps.
    \item \textbf{Rectified Flow}: A laser beam shines directly from your bedroom door to the kitchen refrigerator. You simply walk straight along the laser beam in a direct line. If the beam is clear, you can sprint from start to finish in a single second!
\end{itemize}

\section{Core Mathematical Equations and Definitions}

\subsection{Linear Probability Interpolation Path}
{\boldmath
\begin{equation}
\mathbf{x}_t = (1 - t)\mathbf{x}_0 + t \mathbf{x}_1 = \mathbf{x}_0 + t(\mathbf{x}_1 - \mathbf{x}_0)
\end{equation}
}
\begin{itemize}
    \item \textbf{Formal Definition}: The convex combination mapping base noise $\mathbf{x}_0$ at $t = 0$ to target observation $\mathbf{x}_1$ at $t = 1$.
    \item \textbf{What It Means}: At any time $t \in [0, 1]$, the state is an exact weighted blend of noise and image.
    \item \textbf{Why It Is Important}: Replaces complex trigonometric or variance-preserving schedules with simple linear coordinates, eliminating schedule hyperparameter tuning.
\end{itemize}

\subsection{Canonical Straight-Line Velocity Field}
{\boldmath
\begin{equation}
\mathbf{u}_t(\mathbf{x}_t \mid \mathbf{x}_0, \mathbf{x}_1) = \mathbf{x}_1 - \mathbf{x}_0
\end{equation}
}
\begin{itemize}
    \item \textbf{Formal Definition}: The time derivative $\frac{\mathrm{d}}{\mathrm{d}t}\mathbf{x}_t$ of the straight-line probability trajectory.
    \item \textbf{What It Means}: The velocity depends only on the endpoints and is constant across time for any conditioned trajectory pair $(\mathbf{x}_0, \mathbf{x}_1)$.
    \item \textbf{Why It Is Important}: Provides a closed-form, noise-free regression target during neural network training without requiring backpropagation through ODE solvers.
\end{itemize}

\subsection{Simulation-Free Conditional Flow Matching Loss}
{\boldmath
\begin{equation}
\mathcal{L}_{\text{CFM}}(\theta) = \frac{1}{D} \sum_{i=1}^D \left( \mathbf{v}_\theta(\mathbf{x}_t, t)[i] - (\mathbf{x}_1[i] - \mathbf{x}_0[i]) \right)^2
\end{equation}
}
\begin{itemize}
    \item \textbf{Formal Definition}: The mean squared error regression between the parameterized neural network $\mathbf{v}_\theta$ and the true conditional velocity.
    \item \textbf{What It Means}: Trains the neural network to output the average velocity vector that transports points from noisy regions to clean data.
    \item \textbf{Why It Is Important}: Provably yields the exact same expected gradients as the intractable Marginal Flow Matching loss.
\end{itemize}

\subsection{One-Step Straight Euler Projection}
{\boldmath
\begin{equation}
\hat{\mathbf{x}}_1 = \mathbf{x}_t + (1 - t) \mathbf{v}_\theta(\mathbf{x}_t, t)
\end{equation}
}
\begin{itemize}
    \item \textbf{Formal Definition}: The single-step linear extrapolation from intermediate state $\mathbf{x}_t$ to the destination data boundary $t = 1$.
    \item \textbf{What It Means}: Follows the predicted constant velocity ray directly to the clean image plane.
    \item \textbf{Why It Is Important}: Forms the foundational basis of modern ultra-fast distillation and 1-step Rectified Flow generative pipelines.
\end{itemize}

\section{Rigorous Mathematical Analysis Checklist}

\subsection{Notation and Symbol Semantics}
\begin{itemize}
    \item $\mathbf{x}_0 \in \mathbb{R}^D$: Base prior sample drawn from $p_0 = \mathcal{N}(\mathbf{0}, \mathbf{I})$.
    \item $\mathbf{x}_1 \in \mathbb{R}^D$: Data manifold sample drawn from $p_1$.
    \item $t \in [0, 1]$: Normalized flow time parameter.
    \item $\mathbf{x}_t \in \mathbb{R}^D$: Interpolated state vector at time $t$.
    \item $\mathbf{u}_t \in \mathbb{R}^D$: Conditional ground-truth velocity $\mathbf{x}_1 - \mathbf{x}_0$.
    \item $\mathbf{v}_\theta(\mathbf{x}_t, t) \in \mathbb{R}^D$: Neural network velocity prediction.
\end{itemize}

\subsection{Epistemic Status Table}
\begin{table}[h]
\centering
\small
\begin{tabular}{@{}llll@{}}
\toprule
\textbf{Quantity} & \textbf{Symbol} & \textbf{Epistemic Status} & \textbf{Operational Role} \\
\midrule
Prior Noise & $\mathbf{x}_0$ & Stochastically Sampled & Starting distribution support \\
Target Data & $\mathbf{x}_1$ & Empirically Sampled & Destination manifold support \\
Interpolated State & $\mathbf{x}_t$ & Exact Linear Combination & Intermediate coordinate along path \\
Target Velocity & $\mathbf{u}_t$ & Analytically Exact Derivative & Target regression vector \\
Learned Velocity & $\mathbf{v}_\theta$ & Deep Function Approximation & Marginal transport vector field \\
\bottomrule
\end{tabular}
\end{table}

\subsection{Computational Dependency Graph}
{\centering
$\{\mathbf{x}_0, \mathbf{x}_1, t\} \longrightarrow \{\mathbf{x}_t, \mathbf{u}_t\} \longrightarrow \mathbf{v}_\theta(\mathbf{x}_t, t) \longrightarrow \mathcal{L}_{\text{CFM}} \longrightarrow \hat{\mathbf{x}}_1$
\par}

\subsection{Dimension and Coordinate Consistency}
All states $\mathbf{x}_0, \mathbf{x}_1, \mathbf{x}_t, \mathbf{u}_t, \mathbf{v}_\theta, \hat{\mathbf{x}}_1 \in \mathbb{R}^D$. Time $t \in [0, 1]$ is dimensionless.

\subsection{Asymptotic Regimes}
\begin{itemize}
    \item \textbf{At $t = 0$}: $\mathbf{x}_t = \mathbf{x}_0 \sim \mathcal{N}(\mathbf{0}, \mathbf{I})$; pure prior distribution.
    \item \textbf{At $t = 1$}: $\mathbf{x}_t = \mathbf{x}_1 \sim p_1$; pure data distribution.
    \item \textbf{Reflow limit ($k \to \infty$)}: Trajectories become perfectly straight non-crossing rays; ODE truncation error vanishes.
\end{itemize}

\subsection{Boundary Constraints and Invariants}
\begin{itemize}
    \item \textbf{Endpoint Interpolation Invariant}: $\mathbf{x}_t$ strictly interpolates between $\mathbf{x}_0$ and $\mathbf{x}_1$ with convex weights summing to unity.
\end{itemize}

\subsection{Structural Isomorphisms}
Rectified Flow is isomorphic to the dynamic Benamou-Brenier formulation of Optimal Transport under the quadratic Wasserstein metric $W_2(p_0, p_1)$.

\section{Theoretical Theorems and Formal Proofs}

\begin{theorem}[Equivalence of Conditional and Marginal Flow Matching Gradients]
Let $p_t(\mathbf{x}) = \int p_t(\mathbf{x} \mid \mathbf{x}_1) p_1(\mathbf{x}_1) \mathrm{d}\mathbf{x}_1$ denote the marginal probability path, and let $\mathbf{v}_t(\mathbf{x})$ denote the marginal vector field generating $p_t$. Then the gradient of the intractable Marginal Flow Matching loss $\mathcal{L}_{\text{FM}}(\theta)$ with respect to model parameters $\theta$ is identically equal to the gradient of the tractable Conditional Flow Matching loss $\mathcal{L}_{\text{CFM}}(\theta)$:
{\boldmath
\begin{equation}
\nabla_\theta \mathcal{L}_{\text{FM}}(\theta) = \nabla_\theta \mathcal{L}_{\text{CFM}}(\theta)
\end{equation}
}
\end{theorem}

\begin{proof}
The Marginal Flow Matching objective is defined as:
{\boldmath
\begin{equation}
\mathcal{L}_{\text{FM}}(\theta) = \mathbb{E}_{t, \mathbf{x} \sim p_t(\mathbf{x})}\left[ \left\| \mathbf{v}_\theta(\mathbf{x}, t) - \mathbf{v}_t(\mathbf{x}) \right\|_2^2 \right]
\end{equation}
}
Expanding the squared $L_2$ norm:
{\boldmath
\begin{equation}
\left\| \mathbf{v}_\theta(\mathbf{x}, t) - \mathbf{v}_t(\mathbf{x}) \right\|_2^2 = \|\mathbf{v}_\theta(\mathbf{x}, t)\|_2^2 - 2 \langle \mathbf{v}_\theta(\mathbf{x}, t), \mathbf{v}_t(\mathbf{x}) \rangle + \|\mathbf{v}_t(\mathbf{x})\|_2^2
\end{equation}
}
Since $\|\mathbf{v}_t(\mathbf{x})\|_2^2$ is independent of $\theta$, taking the gradient with respect to $\theta$:
{\boldmath
\begin{equation}
\nabla_\theta \mathcal{L}_{\text{FM}}(\theta) = \mathbb{E}_{t, \mathbf{x} \sim p_t(\mathbf{x})}\left[ 2 \nabla_\theta \mathbf{v}_\theta(\mathbf{x}, t) \left( \mathbf{v}_\theta(\mathbf{x}, t) - \mathbf{v}_t(\mathbf{x}) \right) \right]
\end{equation}
}
By the definition of the marginal vector field via conditional expectation:
{\boldmath
\begin{equation}
\mathbf{v}_t(\mathbf{x}) = \mathbb{E}_{\mathbf{x}_1 \sim p_1}\left[ \mathbf{u}_t(\mathbf{x} \mid \mathbf{x}_1) \mid \mathbf{x}_t = \mathbf{x} \right] = \int \mathbf{u}_t(\mathbf{x} \mid \mathbf{x}_1) \frac{p_t(\mathbf{x} \mid \mathbf{x}_1) p_1(\mathbf{x}_1)}{p_t(\mathbf{x})} \mathrm{d}\mathbf{x}_1
\end{equation}
}
Substituting this conditional expectation into the cross-term:
{\boldmath
\begin{align}
\mathbb{E}_{\mathbf{x} \sim p_t}\left[ \langle \nabla_\theta \mathbf{v}_\theta(\mathbf{x}, t), \mathbf{v}_t(\mathbf{x}) \rangle \right] &= \int p_t(\mathbf{x}) \langle \nabla_\theta \mathbf{v}_\theta(\mathbf{x}, t), \int \mathbf{u}_t(\mathbf{x} \mid \mathbf{x}_1) \frac{p_t(\mathbf{x} \mid \mathbf{x}_1) p_1(\mathbf{x}_1)}{p_t(\mathbf{x})} \mathrm{d}\mathbf{x}_1 \rangle \mathrm{d}\mathbf{x} \\
&= \iint p_1(\mathbf{x}_1) p_t(\mathbf{x} \mid \mathbf{x}_1) \langle \nabla_\theta \mathbf{v}_\theta(\mathbf{x}, t), \mathbf{u}_t(\mathbf{x} \mid \mathbf{x}_1) \rangle \mathrm{d}\mathbf{x}_1 \mathrm{d}\mathbf{x} \\
&= \mathbb{E}_{\mathbf{x}_1 \sim p_1, \mathbf{x} \sim p_t(\cdot \mid \mathbf{x}_1)}\left[ \langle \nabla_\theta \mathbf{v}_\theta(\mathbf{x}, t), \mathbf{u}_t(\mathbf{x} \mid \mathbf{x}_1) \rangle \right]
\end{align}
}
Now consider the Conditional Flow Matching objective:
{\boldmath
\begin{equation}
\mathcal{L}_{\text{CFM}}(\theta) = \mathbb{E}_{t, \mathbf{x}_1 \sim p_1, \mathbf{x} \sim p_t(\cdot \mid \mathbf{x}_1)}\left[ \left\| \mathbf{v}_\theta(\mathbf{x}, t) - \mathbf{u}_t(\mathbf{x} \mid \mathbf{x}_1) \right\|_2^2 \right]
\end{equation}
}
Differentiating $\mathcal{L}_{\text{CFM}}$ with respect to $\theta$ yields the identical expected inner product. Therefore, $\nabla_\theta \mathcal{L}_{\text{FM}}(\theta) = \nabla_\theta \mathcal{L}_{\text{CFM}}(\theta)$ identically.
\end{proof}

\begin{theorem}[Path Straightening and Monotonic Transport Cost Reduction under Reflow]
Let $\mathbf{Z}_1 = \text{ODE}(\mathbf{Z}_0)$ be the terminal points obtained by integrating the 1-Rectified Flow from base samples $\mathbf{Z}_0 \sim p_0$. Then training a 2-Rectified Flow on the paired coupling $(\mathbf{Z}_0, \mathbf{Z}_1)$ strictly decreases the expected transport path length:
{\boldmath
\begin{equation}
\mathbb{E}\left[ \int_0^1 \|\mathbf{v}^{(2)}_t(\mathbf{x}_t)\|_2 \mathrm{d}t \right] \le \mathbb{E}\left[ \int_0^1 \|\mathbf{v}^{(1)}_t(\mathbf{x}_t)\|_2 \mathrm{d}t \right]
\end{equation}
}
with equality holding if and only if the flow trajectories are already straight non-intersecting lines.
\end{theorem}

\begin{proof}
Let $\gamma(t): [0, 1] \to \mathbb{R}^D$ be any continuous trajectory connecting $\mathbf{Z}_0$ to $\mathbf{Z}_1 = \gamma(1)$. The length of the curve is given by $L(\gamma) = \int_0^1 \|\dot{\gamma}(t)\|_2 \mathrm{d}t$.
By Euclidean geometry, the straight line segment $\gamma_{\text{straight}}(t) = (1 - t)\mathbf{Z}_0 + t\mathbf{Z}_1$ is the unique geodesic in $\mathbb{R}^D$ minimizing length between endpoints:
{\boldmath
\begin{equation}
L(\gamma_{\text{straight}}) = \|\mathbf{Z}_1 - \mathbf{Z}_0\|_2 = \left\| \int_0^1 \dot{\gamma}(t) \mathrm{d}t \right\|_2 \le \int_0^1 \|\dot{\gamma}(t)\|_2 \mathrm{d}t = L(\gamma)
\end{equation}
}
The 1-Rectified Flow was trained on independent couplings $(\mathbf{x}_0, \mathbf{x}_1) \sim p_0 \times p_1$, resulting in curved trajectories due to vector field crossing.
Integrating the ODE from $\mathbf{Z}_0$ yields an exact causal coupling $(\mathbf{Z}_0, \mathbf{Z}_1)$ with no crossing paths at endpoints.
Training the 2-Rectified Flow on straight lines connecting $(\mathbf{Z}_0, \mathbf{Z}_1)$ replaces the curved integral paths with straight chord segments.
Taking expectations over the joint distribution:
{\boldmath
\begin{equation}
\mathbb{E}\left[ \|\mathbf{Z}_1 - \mathbf{Z}_0\|_2 \right] \le \mathbb{E}\left[ \int_0^1 \|\mathbf{v}^{(1)}_t(\mathbf{x}_t)\|_2 \mathrm{d}t \right]
\end{equation}
}
Thus, each reflow iteration progressively straightens the vector field, strictly reducing path curvature and ODE discretization truncation error.
\end{proof}

\section{Foundational Architectural Questions and Epistemic Meta-Questions}

\subsection{Architectural Question 1: Flow Matching vs Diffusion Models}
\textbf{Question:} Why has Flow Matching replaced traditional DDPM in modern state-of-the-art models (Stable Diffusion 3, Flux)?\\
\textbf{Answer:} DDPM forces curved trajectory schedules (VP-SDE) that create high curvature near $t = 0$. Flow Matching establishes straight-line linear vector fields between noise and data, drastically reducing numerical ODE stiffness and enabling photorealistic generation in 4 to 8 steps without quality degradation.

\textbf{Meta-Question:} Is Flow Matching fundamentally a different mathematical model from diffusion?\\
\textbf{Answer:} No. Diffusion models can be reparameterized as Continuous Normalizing Flows via the Probability Flow ODE. Flow Matching is a generalized framework for training Continuous Normalizing Flows with arbitrary probability paths, of which diffusion is simply one sub-optimal curved instance.

\subsection{Architectural Question 2: The Reflow Procedure}
\textbf{Question:} What is the purpose of the ``Reflow'' step, and why can't the initial model achieve 1-step generation directly?\\
\textbf{Answer:} In the initial training (1-Rectified Flow), pairs $(\mathbf{x}_0, \mathbf{x}_1)$ are drawn independently, causing straight-line paths to cross in high dimensions. The neural network learns the conditional expectation of these intersecting rays, resulting in curved trajectories. Reflow pairs $\mathbf{x}_0$ with its deterministic ODE output $\mathbf{x}_1$, uncrossing trajectories and straightening the vector field.

\textbf{Meta-Question:} How many reflow iterations are necessary before trajectories become completely straight?\\
\textbf{Answer:} In practice, a single Reflow iteration (2-Rectified Flow) uncrosses over 95\% of trajectory curvature, allowing distillation or Euler solvers to generate high-fidelity samples in 1 to 2 steps.

\section{Algorithmic Implementation Specification}

\begin{algorithm}[H]
\caption{Flow Matching Interpolation and Straight-Line Objective Evaluation}
\begin{algorithmic}[1]
\Require Noise sample $\mathbf{x}_0 \in \mathbb{R}^D$, data sample $\mathbf{x}_1 \in \mathbb{R}^D$, predicted velocity $\mathbf{v}_\theta \in \mathbb{R}^D$, timestep $t \in [0, 1]$
\Ensure Interpolated state $\mathbf{x}_t \in \mathbb{R}^D$, target velocity $\mathbf{u}_t \in \mathbb{R}^D$, CFM loss $\mathcal{L}_{\text{CFM}}$, 1-step sample $\hat{\mathbf{x}}_1 \in \mathbb{R}^D$
\State Initialize empty arrays $\mathbf{x}_t, \mathbf{u}_t, \hat{\mathbf{x}}_1$ of length $D$
\State $\text{loss\_sum} \gets 0.0$
\For{$i = 0$ \textbf{to} $D - 1$}
    \State $\mathbf{x}_t[i] \gets (1 - t) \cdot \mathbf{x}_0[i] + t \cdot \mathbf{x}_1[i]$
    \State $\mathbf{u}_t[i] \gets \mathbf{x}_1[i] - \mathbf{x}_0[i]$
    \State $diff \gets \mathbf{v}_\theta[i] - \mathbf{u}_t[i]$
    \State $\text{loss\_sum} \gets \text{loss\_sum} + diff \cdot diff$
    \State $\hat{\mathbf{x}}_1[i] \gets \mathbf{x}_t[i] + (1 - t) \cdot \mathbf{v}_\theta[i]$
\EndFor
\State $\mathcal{L}_{\text{CFM}} \gets \text{loss\_sum} / D$
\State \Return $(\mathbf{x}_t, \mathbf{u}_t, \mathcal{L}_{\text{CFM}}, \hat{\mathbf{x}}_1)$
\end{algorithmic}
\end{algorithm}

\section{Big-$\Theta$ Complexity Analysis}

\begin{itemize}
    \item \textbf{Interpolation and Loss Time Complexity}: $\Theta(D)$ vector arithmetic operations per training sample.
    \item \textbf{Total Sampling Generation Time Complexity}: For an $S$-step solver, total runtime is $\Theta(S \cdot (D + \mathcal{C}_{\text{NN}}))$, where $S \in [1, 4]$ for rectified flows.
    \item \textbf{Space Complexity}: $\Theta(D)$ working buffers to maintain spatial states.
    \item \textbf{Work \& Depth}: Work is $\mathcal{W} = \Theta(S \cdot D)$; parallel depth across spatial coordinates is $\mathcal{D} = \Theta(S \cdot \log D)$.
\end{itemize}

\section{Single Canonical Repository Reference}

The complete deterministic scalar implementation, verification suite, and unit test benchmarks are published at:
\begin{center}
\url{https://github.com/Chief-Strategist-J/llm-obs-infra}
\end{center}

\end{document}
